\documentclass[11pt,a4paper]{article}
\usepackage[utf8]{inputenc}
\usepackage{geometry}
\geometry{margin=0.75in}

\begin{document}

\begin{center}
    {\LARGE \textbf{Morgan Chen}} \\
    \vspace{2pt}
    \textbf{Applied AI \& LLM Systems Engineer} \\
    New York, NY \textbar\ morgan.chen.synthetic@example.com \textbar\ github.com/morganchen-synthetic
\end{center}

\section*{Summary}
Applied AI Engineer specializing in LLM deployment, fine-tuning, retrieval-augmented generation (RAG), and high-throughput inference optimization on GPU clusters.

\section*{Technical Skills}
\textbf{AI \& ML Frameworks:} PyTorch, Hugging Face Transformers, vLLM, TensorRT-LLM, LangGraph, LlamaIndex, Triton Inference Server \\
\textbf{Languages:} Python, C++, CUDA, SQL \\
\textbf{Infrastructure \& Cloud:} Kubernetes, Ray, AWS SageMaker, Docker, Weights \& Biases

\section*{Professional Experience}
\textbf{Senior Applied AI Engineer} \hfill 2022 -- Present \\
\textit{OmniAI Technologies} \hfill New York, NY
\begin{itemize}
    \item Built enterprise RAG pipeline serving 2M daily semantic queries using LangGraph, Milvus, and hybrid sparse-dense retrieval with sub-50ms latency.
    \item Optimized vLLM inference cluster throughput by 3.2x using FP8 quantization and continuous batching on NVIDIA H100 clusters.
    \item Fine-tuned 8B and 70B parameter open-source models using LoRA/QLoRA for domain-specific code generation with 91\% benchmark accuracy.
\end{itemize}

\section*{Education}
\textbf{M.S. in Machine Learning} \hfill 2020 -- 2022 \\
Columbia University, New York, NY

\end{document}
